\documentclass[a4paper,11pt]{article}
\usepackage{exam-style-341}
\examinehead{341 农业综合知识三 · 数据库分章练习}{第03章\quad SQL 语言}
\begin{document}

\unitcover{341 农业综合知识三 · 数据库分章练习}{第03章\quad SQL 语言}{练习（共 100 分）}{约 60 分钟}{SQL 语言的特点与组成；数据定义（CREATE、ALTER、DROP）与约束（主键、外键、CHECK、NOT NULL）；数据操纵（INSERT、UPDATE、DELETE）；数据查询（单表查询、连接查询、嵌套查询、EXISTS、GROUP BY、HAVING、ORDER BY、聚合函数）；视图的定义、作用及与基本表的区别；索引；权限管理 GRANT 与 REVOKE。}

\begin{knowmap}
\textbf{SQL 的组成：}数据定义 DDL（CREATE、ALTER、DROP）· 数据操纵 DML（INSERT、UPDATE、DELETE）· 数据查询 DQL（SELECT）· 数据控制 DCL（GRANT、REVOKE）· 事务控制（COMMIT、ROLLBACK）\\[2pt]
\textbf{SELECT 的书写顺序：}SELECT ... FROM ... WHERE ... GROUP BY ... HAVING ... ORDER BY ...；执行顺序为 FROM $\rightarrow$ WHERE $\rightarrow$ GROUP BY $\rightarrow$ HAVING $\rightarrow$ SELECT $\rightarrow$ ORDER BY\\[2pt]
\textbf{查询层次：}单表查询 $\rightarrow$ 连接查询（等值连接、自然连接、自身连接、外连接）$\rightarrow$ 嵌套查询（IN、比较运算符、EXISTS、相关子查询）$\rightarrow$ 集合查询\\[2pt]
\textbf{分组与聚合：}COUNT、SUM、AVG、MAX、MIN 作用于组；WHERE 在分组前筛选元组，HAVING 在分组后筛选组，组函数只能出现在 SELECT 列表和 HAVING 子句中\\[2pt]
\textbf{视图：}由基本表或其他视图导出的虚表，数据库中只保存视图定义，不保存视图数据；视图可以简化查询、提供逻辑独立性和安全保护，但并非所有视图都可更新\\[2pt]
\textbf{数据控制：}GRANT 授权、REVOKE 收回权限；权限可授予用户或角色，并可用 WITH GRANT OPTION 允许转授权限
\end{knowmap}

\begin{keyterms}
\textbf{SQL}：结构化查询语言，是关系数据库的标准语言，具有综合统一、高度非过程化、面向集合、可嵌入等特点。\\
\textbf{DDL}：数据定义语言，用于定义、修改和删除数据库对象，如 CREATE、ALTER、DROP。\\
\textbf{DML}：数据操纵语言，用于插入、更新和删除数据，如 INSERT、UPDATE、DELETE。\\
\textbf{DQL}：数据查询语言，主要指 SELECT 语句及其各种子句。\\
\textbf{DCL}：数据控制语言，用于权限管理，如 GRANT、REVOKE。\\
\textbf{主键约束}：用 PRIMARY KEY 定义，要求主码列值唯一且非空，一个表只能有一个主键。\\
\textbf{外键约束}：用 FOREIGN KEY ... REFERENCES 定义，要求外码值等于被参照表中某个主键值或为空值。\\
\textbf{空值}：NULL，表示未知或不适用；判断空值必须用 IS NULL 或 IS NOT NULL，不能用 = NULL。\\
\textbf{连接查询}：同时涉及两个或两个以上表的查询，通过连接条件把相关元组拼接起来。\\
\textbf{嵌套查询}：将一个 SELECT 语句嵌套在另一个 SELECT 语句的 WHERE 或 HAVING 子句中的查询。\\
\textbf{相关子查询}：子查询的执行依赖于外层查询的当前元组，外层每取一个元组，子查询就执行一次。\\
\textbf{聚合函数}：对一组值进行计算并返回单个值的函数，如 COUNT、SUM、AVG、MAX、MIN。\\
\textbf{视图}：从一个或几个基本表（或其他视图）导出的虚表，数据库中只存放视图的定义而不存放视图对应的数据。\\
\textbf{索引}：为加快数据检索而建立的数据结构，可以建立在一个或多个属性列上，分为普通索引、唯一索引等。
\end{keyterms}

\begin{contrast}{WHERE / GROUP BY / HAVING / ORDER BY（每行依次为：作用·能否使用聚合函数·执行顺序与典型示例）}
\citem{WHERE}{筛选满足条件的元组}{不能使用聚合函数}{FROM 之后、GROUP BY 之前执行，如 WHERE Price > 10}
\citem{GROUP BY}{按指定列的值对元组分组}{不能使用聚合函数}{WHERE 之后、HAVING 之前执行，如 GROUP BY Cid}
\citem{HAVING}{筛选满足条件的组}{可以使用聚合函数}{GROUP BY 之后、SELECT 之前执行，如 HAVING SUM(Qty) > 100}
\citem{ORDER BY}{对最终结果排序}{可以对结果列使用}{SELECT 之后、最后执行，如 ORDER BY Total DESC}
\end{contrast}

\begin{contrast}{基本表 / 视图（每行依次为：数据存储·可更新性·主要作用）}
\citem{基本表}{实际存储数据}{一般可直接更新}{数据库中数据的实际组织单位}
\citem{视图}{只存定义，不存数据（虚表）}{受定义限制，并非都可更新}{简化查询、逻辑独立性、安全保护、多角度看待数据}
\end{contrast}

\begin{contrast}{DELETE / DROP / TRUNCATE（每行依次为：删除对象·是否删除表定义·能否有条件及事务回滚）}
\citem{DELETE}{表中的元组（行）}{不删除表定义}{可带 WHERE 条件，属于 DML，可回滚}
\citem{DROP}{整个表（定义和数据）}{删除表定义}{不能带条件，属于 DDL，一般不能回滚}
\citem{TRUNCATE}{表中的全部元组}{保留表定义}{不能带条件，属于 DDL，一般不能回滚，速度较快}
\end{contrast}

\begin{contrast}{连接查询 / 嵌套子查询 / EXISTS 子查询（每行依次为：基本结构·返回结果·适用场合）}
\citem{连接查询}{FROM 中列出多个表并用连接条件连接}{一张包含多表属性的结果表}{需要同时显示多个表的属性时}
\citem{嵌套子查询}{WHERE 中嵌套 SELECT，常用 IN 或比较运算符}{子查询返回集合或单值}{子查询结果作为外层查询的筛选条件时}
\citem{EXISTS 子查询}{WHERE EXISTS（子查询）}{只判断子查询是否为空，返回真或假}{相关子查询中判断“是否存在满足条件的记录”时}
\end{contrast}

\qtypename{一、单项选择题}{10 题，每题 2 分，共 20 分}
\q SQL 语言中，修改基本表结构的语句是（\quad）
\choicesline{A. MODIFY TABLE}{B. ALTER TABLE}{C. UPDATE TABLE}{D. CHANGE TABLE}

\q 在 SELECT 语句中，若要消除查询结果中的重复元组，应使用（\quad）
\choicesline{A. DISTINCT}{B. UNIQUE}{C. GROUP BY}{D. ORDER BY}

\q 在 SQL 中，条件“姓名 LIKE '李\%'”的含义是（\quad）
\choices{姓名以“李”开头}{姓名以“李”结尾}{姓名中包含“李”}{姓名恰好为“李”}

\q 下列聚合函数中，用于求某列平均值的是（\quad）
\choicesline{A. SUM}{B. COUNT}{C. AVG}{D. MAX}

\q 下列关于 WHERE 与 HAVING 的说法中，正确的是（\quad）
\choices{WHERE 子句可以使用聚合函数}{HAVING 子句用于对分组后的结果进行筛选}{HAVING 子句在 GROUP BY 之前执行}{WHERE 子句只能与 GROUP BY 一起使用}

\q 视图是一个“虚表”，其含义是（\quad）
\choices{视图中只存放视图的定义，不存放视图对应的数据}{视图中存放的数据是基本表数据的副本}{视图不能用于查询}{视图必须由多个基本表导出}

\q 若要删除基本表及其全部数据，应使用的语句是（\quad）
\choicesline{A. DELETE FROM 表名}{B. DROP TABLE 表名}{C. TRUNCATE 表名}{D. REMOVE TABLE 表名}

\q 在相关子查询中，子查询的执行次数通常取决于（\quad）
\choices{外层查询结果中的元组个数}{子查询中表的个数}{内层查询结果中的元组个数}{SELECT 子句中列的个数}

\q 语句 GRANT SELECT ON SC TO U1 的作用是（\quad）
\choices{把对表 SC 的查询权限授予用户 U1}{把用户 U1 的查询权限收回}{创建用户 U1}{把表 SC 的所有权转给 U1}

\q 设 SC 表中有 100 个元组，其中 Grade 列有 10 个空值，则 COUNT(*) 和 COUNT(Grade) 的结果分别是（\quad）
\choicesline{A. 100，90}{B. 90，100}{C. 100，100}{D. 90，90}

\qtypename{二、多项选择题}{5 题，每题 4 分，共 20 分；每题有两个或两个以上正确项}
\q SQL 语言的主要特点包括（\quad）
\choices{综合统一，集数据定义、操纵、控制于一体}{高度非过程化}{面向集合的操作方式}{只能作为独立语言使用，不能嵌入其他语言}

\q 下列属于 SQL 聚合函数的有（\quad）
\choicesline{A. COUNT}{B. SUM}{C. AVG}{D. LENGTH}

\q 关于视图，下列说法正确的有（\quad）
\choices{视图是从一个或几个基本表导出的虚表}{视图能够简化用户的操作}{视图能够对机密数据提供安全保护}{所有视图都可以直接更新}

\q 使用 GROUP BY 子句时，下列做法正确的有（\quad）
\choices{SELECT 列表中只能出现分组属性和聚合函数}{HAVING 子句中可以使用聚合函数}{WHERE 子句中可以出现聚合函数}{GROUP BY 子句可以按多个属性列分组}

\q 关于索引，下列说法正确的有（\quad）
\choices{建立索引可以加快数据查询的速度}{建立索引会降低插入和更新数据的速度}{索引可以建立在一个或多个属性列上}{建立索引会改变基本表的逻辑结构}

\qtypename{三、判断题}{5 题，每题 2 分，共 10 分；正确的画“√”，错误的画“×”并写出改正后的表述}
\q SQL 是一种过程化语言，使用 SQL 时必须指出具体的数据存取路径。\quad（\quad）\quad 改正：\fixblank

\q WHERE 子句中可以使用 COUNT、SUM 等聚合函数。\quad（\quad）\quad 改正：\fixblank

\q 视图是一个虚表，数据库中只存放视图的定义，视图的数据仍存放在基本表中。\quad（\quad）\quad 改正：\fixblank

\q DROP TABLE 删除基本表后，其上建立的视图会失效。\quad（\quad）\quad 改正：\fixblank

\q COUNT(列名) 会统计该列中的空值。\quad（\quad）\quad 改正：\fixblank

\qtypename{四、填空题}{5 题，每题 2 分，共 10 分}
\q SQL 的功能可分为数据定义、\fillblank{} 和数据控制，其中数据定义使用的语句主要有 CREATE、ALTER 和 \fillblank{}。

\q SQL 中用于模糊匹配的通配符中，“\%”表示 \fillblank{}，“\_”表示 \fillblank{}。

\q SQL 中常用的聚合函数有 COUNT、SUM、AVG、\fillblank{} 和 \fillblank{}。

\q 创建视图使用 \fillblank{} 语句；把权限授予用户使用 \fillblank{} 语句，收回权限使用 REVOKE 语句。

\q SELECT 语句中子句的执行顺序依次为 FROM、\fillblank{}、GROUP BY、\fillblank{}、SELECT、ORDER BY。

\qtypename{五、简答题}{2 题，每题 8 分，共 16 分}
\q 什么是视图？视图与基本表有什么区别？使用视图有哪些优点？\anslines{2.2cm}

\q 试说明 WHERE、GROUP BY、HAVING 和 ORDER BY 四个子句的作用和执行顺序，并说明为什么聚合函数不能出现在 WHERE 子句中。\anslines{2.2cm}

\qtypename{六、综合题}{2 题，每题 12 分，共 24 分}
\q 某农产品销售数据库中有如下五个基本表（带下划线的属性为主码）：
\begin{center}
\small
\begin{tabular}{@{}lll@{}}
\toprule
关系名 & 属性（含义） & 说明 \\
\midrule
Farmer & \underline{FarmerId}, Fname, Phone & 农户编号、姓名、电话 \\
Product & \underline{Pid}, Pname, Price, FarmerId & 产品编号、名称、单价、所属农户 \\
Customer & \underline{Cid}, Cname, City & 客户编号、姓名、所在城市 \\
Orders & \underline{Oid}, Odate, Cid & 订单号、订单日期、下单客户 \\
OrderItem & \underline{Oid}, \underline{Pid}, Qty & 订单号、产品编号、数量 \\
\bottomrule
\end{tabular}
\end{center}

其中：Product 的 FarmerId 参照 Farmer 的 FarmerId；Orders 的 Cid 参照 Customer 的 Cid；OrderItem 的 Oid 参照 Orders 的 Oid，Pid 参照 Product 的 Pid。

请完成下列各题：
\begin{enumerate}[label=（\arabic*）, leftmargin=2.4em, itemsep=1pt]
  \item 用 CREATE TABLE 语句创建 OrderItem 表，要求定义复合主码 (Oid, Pid)，定义两个外码，并约束 Qty 非空且大于 0。（3 分）
  \item 用 ALTER TABLE 语句为 Farmer 表增加一列 Email，类型为 VARCHAR(50)。（2 分）
  \item 向 OrderItem 表中插入一条记录：订单号 'O001'、产品编号 'P01'、数量 10。（2 分）
  \item 将订单 'O001' 中产品 'P01' 的数量增加 5。（2 分）
  \item 删除订单 'O001' 及其全部订单明细，并说明删除顺序为什么不能颠倒。（3 分）
\end{enumerate}
\anslines{4.5cm}

\q 沿用第 28 题的五个基本表，完成下列查询：
\begin{enumerate}[label=（\arabic*）, leftmargin=2.4em, itemsep=1pt]
  \item 查询单价高于 10 元的产品名称及其所属农户的姓名。（3 分）
  \item 查询没有下过任何订单的客户姓名。（3 分）
  \item 查询订单总金额超过 1000 元的订单号及订单总金额，结果按订单号升序排列。（3 分）
  \item 创建视图 V\_ProductSales，统计每种产品的产品编号、产品名称、销售总量和销售总金额，并查询销售总金额超过 500 元的产品。（3 分）
\end{enumerate}
\anslines{4.5cm}

\end{document}
